\documentclass[11pt,letterpaper]{article}
\usepackage[margin=1in]{geometry}
\usepackage{amsmath,amssymb,amsthm}
\usepackage{graphicx}
\usepackage{hyperref}
\usepackage{booktabs}
\usepackage{enumitem}
\usepackage{microtype}
\usepackage{xcolor}

\hypersetup{
  colorlinks=true, linkcolor=blue!60!black, citecolor=blue!60!black,
  urlcolor=blue!60!black,
}

\title{\Large\textbf{Distance-3 Steane code repeated-round logical-qubit experiment on \texttt{ibm\_kingston}}\\
\large{$P(\hat L = 0) = 0.6235 \pm 0.0076$ at 4 syndrome rounds, $16.3\sigma$ above random, on a Heron-r2 processor}}

\author{Christian Knopp \\
  \texttt{cknopp@gmail.com} \\
  \href{https://github.com/Zynerji/systrophe}{\texttt{github.com/Zynerji/systrophe}}}

\date{2026-05-12 \quad / \quad \Systrophe{}\ v0.19.x post-tag}

\providecommand{\Systrophe}{Systroph\=e}

\begin{document}
\maketitle

\begin{abstract}
We report a real distance-3 Steane \texttt{[[7,1,3]]} code repeated-round logical-qubit experiment on IBM Quantum's \texttt{ibm\_kingston} 156-qubit Heron-r2 processor, paired with a wall-clock-matched bare-qubit baseline. The experiment prepares $|0_L\rangle$ via the canonical CSS encoder on 7 data qubits, performs $n_{\mathrm{rounds}} \in \{1, 2, 4, 8\}$ rounds of mid-circuit X- and Z-syndrome extraction with 6 ancilla qubits, and decodes the data-measurement-derived syndrome via lookup table. \textbf{The logical-Z=0 rate decays smoothly from $0.81$ at $n_{\mathrm{rounds}} = 1$ to $0.53$ at $n_{\mathrm{rounds}} = 8$, but it is $1.7$–$9.1$ percentage points \emph{below} the corresponding bare-qubit baseline at every round count}, demonstrating that the $d = 3$ Steane code at current Heron-r2 gate error rates ($\mathrm{cz}_{\mathrm{err}} \approx 2 \times 10^{-3}$) is \emph{sub-threshold}: encoding + syndrome overhead introduces more errors than the $d = 3$ code can correct. This is the standard expected behaviour and the canonical reason why Google Quantum AI's surface-code demonstrations required $d \ge 5$ to reach break-even \cite{willow_precursor,willow}. Our experiment establishes a clean Heron-r2 baseline against which future higher-distance \Systrophe{} hardware experiments will be compared, and supersedes our earlier 156-qubit GHZ majority-vote demonstration which we have retracted as not-QEC.
\end{abstract}

\section{Introduction and scope}\label{sec:intro}

Quantum error correction (QEC) is the defining engineering challenge for fault-tolerant quantum computation. The Clay-level mathematical results that motivated Shor and Steane (the existence of $[[n,k,d]]$ stabilizer codes with $d \ge 3$ on $n = O(\mathrm{poly}\,k)$ physical qubits) have been demonstrated in increasingly large hardware experiments: notably Google Quantum AI's Willow distance-7 surface code below threshold (December 2024) \cite{willow}, IBM's bivariate-bicycle \texttt{[[144,12,12]]} construction on heavy-hex topologies \cite{bb_paper}, and earlier distance-3 demonstrations on superconducting and ion-trap platforms.

This paper reports a more modest experiment: a real distance-3 Steane \texttt{[[7,1,3]]} CSS code with 4 rounds of mid-circuit syndrome extraction on a 156-qubit Heron-r2 processor (\texttt{ibm\_kingston}). The experiment uses 13 physical qubits total (7 data, 6 ancilla) and runs as a single SamplerV2 job. The aims:
\begin{enumerate}[leftmargin=*,itemsep=0pt]
\item Verify that a textbook CSS code can be executed end-to-end on Heron-r2 with the IBM Qiskit Runtime defaults (dynamical decoupling on, XpXm sequence).
\item Measure the residual logical Z fidelity after 4 rounds at the ISA-transpiled depth ($\approx 556$).
\item Establish a clean baseline against which future higher-distance experiments in the \Systrophe{} framework will be compared.
\item Retract our earlier \texttt{156-qubit GHZ majority-vote} result as a QEC claim. That experiment was a Z-basis GHZ measurement followed by classical post-processing; the repetition code protects only against bit-flip errors and the 99.1\% recovery rate is binomial-CDF arithmetic, not error correction.
\end{enumerate}

\section{Steane \texttt{[[7,1,3]]} code: encoder and decoder}\label{sec:code}

The Steane code is a self-dual CSS code on 7 physical qubits with $k=1$ logical qubit and distance $d=3$, correcting any single-qubit error. We use a parity-check matrix
\[
H_{\Systrophe} \;=\; \begin{pmatrix}
1 & 1 & 0 & 1 & 1 & 0 & 0 \\
1 & 0 & 1 & 1 & 0 & 1 & 0 \\
0 & 1 & 1 & 1 & 0 & 0 & 1
\end{pmatrix}
\]
which corresponds to the canonical |0_L> encoder
\[
\begin{aligned}
&H \!\cdot q_4,\ H \!\cdot q_5,\ H \!\cdot q_6 \\
&\mathrm{CNOT}(q_4 \!\to\! q_0),\ \mathrm{CNOT}(q_4 \!\to\! q_1),\ \mathrm{CNOT}(q_4 \!\to\! q_3) \\
&\mathrm{CNOT}(q_5 \!\to\! q_0),\ \mathrm{CNOT}(q_5 \!\to\! q_2),\ \mathrm{CNOT}(q_5 \!\to\! q_3) \\
&\mathrm{CNOT}(q_6 \!\to\! q_1),\ \mathrm{CNOT}(q_6 \!\to\! q_2),\ \mathrm{CNOT}(q_6 \!\to\! q_3)
\end{aligned}
\]
which prepares the 8-term equal-superposition
\[
|0_L\rangle \;=\; \tfrac{1}{\sqrt 8}\!\bigl(|0000000\rangle + |1101100\rangle + |1010101\rangle + |0111001\rangle + |0110110\rangle + |1011011\rangle + |1100011\rangle + |0001111\rangle\bigr).
\]
All 8 codewords have even Hamming weight, so the transversal $Z_L = \prod_q Z_q$ is $+1$ on $|0_L\rangle$.

\subsection{Syndrome extraction}

Each round of syndrome extraction adds 6 ancilla measurements: 3 X-stabilizers and 3 Z-stabilizers. The X-stab ancillas are prepared in $|+\rangle$ via Hadamard, entangled with the data qubits in the stabilizer support via CNOT, mapped back to $|0/1\rangle$ via Hadamard, and measured (with reset afterwards). The Z-stab ancillas are prepared in $|0\rangle$, take CNOTs FROM the data qubits, and are measured. The 6 syndrome bits per round are saved to classical registers for offline decoding.

\subsection{Decoder}

We use a simple lookup-table decoder on the data-measurement-derived syndrome:
\[
s_i \;=\; \bigl(\sum_{q=0}^{6} H_{\Systrophe,\,i,q} \cdot \mathrm{data}_q\bigr) \bmod 2.
\]
If $\mathbf{s} \ne (0,0,0)$, the lookup table identifies the column of $H_{\Systrophe}$ matching $\mathbf{s}$ and flips the corresponding data bit. The logical Z is then computed as the parity of the 7 (corrected) data bits. For $d = 3$ this single-round data-syndrome decoder is sufficient; for higher distances, the proper decoder is minimum-weight perfect matching on the 3D syndrome graph.

\subsection{Noiseless validation}

In noiseless Qiskit-Aer simulation on the same circuit, the decoder returns $P(\hat L = 0) = 1.0$ exactly, for $n_{\mathrm{rounds}} \in \{0, 1, 3\}$, validating the encoder + decoder + syndrome-extraction circuit (Section~\ref{sec:hw}). Each data qubit measures $\approx 0.5$ in $Z$ as expected, since the $|0_L\rangle$ data register is in an equal superposition of 8 codewords.

\section{Hardware experiment on \texttt{ibm\_kingston}}\label{sec:hw}

\subsection{Submission}

The full 4-round circuit (encoding + 4 rounds of syndrome extraction + final data measurement) was submitted to \texttt{ibm\_kingston} via the Qiskit Runtime SamplerV2 with the following options:
\begin{itemize}[leftmargin=*,itemsep=0pt]
\item Backend: \texttt{ibm\_kingston} (156-qubit Heron-r2)
\item Qubits used: 13 (7 data, 6 ancilla)
\item Optimization level: 3
\item ISA-transpiled depth: 556
\item Dynamical decoupling: enabled, XpXm sequence
\item Shots: 4096
\item Job ID: \texttt{d81k1s7oha1c73bkggg0}
\item Submitted: 2026-05-12 (after the IBM Quantum allocation-warning that was issued in the same allocation window did not prevent prior batch-7 jobs from running).
\end{itemize}

\subsection{Result}

The job completed successfully. Decoding the 4096-shot measurement yields:

\begin{center}
\begin{tabular}{lr}
\toprule
Quantity & Value \\
\midrule
Logical Z = 0 rate $P(\hat L = 0)$ & $\bm{0.6235}$ \\
$1\sigma$ shot-noise uncertainty & $\pm 0.0076$ \\
Distance from random ($0.5$) & $+12.35$ pp \\
Distance from random (in $\sigma$) & $\bm{+16.3\,\sigma}$ \\
Distance from noiseless ideal ($1.0$) & $-37.65$ pp \\
Mean physical $Z = 0$ rate per qubit & $0.4998$ \\
\bottomrule
\end{tabular}
\end{center}

The mean physical $Z = 0$ rate is $0.4998$, consistent with each data qubit being in an equal-weight superposition (which is what $|0_L\rangle$ predicts for each individual data qubit). The logical recovery at $62.35\%$ is significantly above the $50\%$ random baseline at $16.3\sigma$, demonstrating that the Steane code is actively preserving logical information beyond what random per-qubit measurement would give.

\subsection{Interpretation}

The result establishes three facts about Heron-r2 Steane-code QEC:

\begin{enumerate}[leftmargin=*,itemsep=0pt]
\item \textbf{The code runs end-to-end.} 4 rounds of mid-circuit syndrome extraction with reset complete on Heron-r2 hardware without errors. The ISA-transpiled circuit at depth 556 stays coherent enough that decoder recovery is meaningful.
\item \textbf{Logical signal is well above random.} A 12-percentage-point excess at 16.3$\sigma$ is unambiguous: the encoding + syndrome + decode pipeline is doing real work.
\item \textbf{We are not below threshold.} A 38-percentage-point deficit below noiseless ideal means decoherence over the 4-round, depth-556 ISA circuit is substantial. A direct \texttt{logical vs physical} comparison at matched circuit depth (i.e., a bare-qubit baseline circuit of the same elapsed time) was not run in this experiment; we cannot yet make a \emph{logical error rate $<$ physical} claim.
\end{enumerate}

\subsection{Comparison with the published QEC state of the art}\label{ssec:sota}

\textbf{Google Quantum AI Willow (Dec 2024)} demonstrated distance-7 surface-code repetition with logical error rate \emph{below} physical for the first time at superconducting scale; their experiment uses real fault-tolerance with much higher distance and many more rounds.

\textbf{IBM bivariate-bicycle codes (2024)} achieved the \texttt{[[144,12,12]]} code on heavy-hex topology with high tolerance to syndrome-measurement errors.

By comparison, our \texttt{[[7,1,3]]} Steane code at $n_{\mathrm{rounds}} = 4$ on Heron-r2:
\begin{itemize}[leftmargin=*,itemsep=0pt]
\item Uses the lowest non-trivial distance ($d = 3$).
\item Runs only 4 rounds.
\item Uses a lookup-table decoder rather than minimum-weight matching.
\item Does not yet include a logical-vs-physical lifetime comparison.
\end{itemize}

\textbf{This experiment is not SOTA QEC.} It is a verified Heron-r2 distance-3 logical-qubit baseline. The honest claim-level (per the framework's 3-tier taxonomy):
\begin{itemize}[leftmargin=*,itemsep=0pt]
\item \emph{Reproducible}: yes, anyone with Qiskit + IBM Quantum access can rerun.
\item \emph{Novel demonstration on Heron-r2}: yes, in the narrow sense that we landed a distance-3 Steane code with 4 mid-circuit syndrome rounds at 4096 shots on \texttt{ibm\_kingston}.
\item \emph{SOTA contribution}: \textbf{no} (does not beat Willow or BB-code published thresholds).
\end{itemize}

\section{Retraction of the 156-qubit GHZ majority-vote result}\label{sec:retraction}

In an earlier press kit (\texttt{launch/viral\_press\_release\_156q\_qec.md}, since superseded), we described a 156-qubit GHZ + majority-vote experiment on \texttt{ibm\_kingston} as "distance-156 repetition-code QEC at 99.1\% logical fidelity." That framing does not survive review. The Z-basis measurement of a 156-qubit GHZ produces a Hamming-weight distribution; majority-voting that distribution is binomial-CDF arithmetic on a noisy classical bit, not error correction. The repetition code protects only against bit-flip errors; the GHZ state has both bit-flip and phase-flip vulnerabilities, and the latter is not detected by the majority vote.

The earlier draft has been retracted (file rewritten as an honest post-mortem); the result is not for distribution as a QEC headline. The present Steane experiment is the real first-step QEC artifact for the \Systrophe{} framework on Heron-r2.

\section{Round sweep + physical baseline}\label{sec:round_sweep}

We additionally ran a paired round sweep with a matched bare-qubit baseline. For each $n_{\mathrm{rounds}} \in \{1, 2, 4, 8\}$, we submitted both:
\begin{enumerate}[leftmargin=*,itemsep=0pt]
\item The full Steane experiment as in Section~\ref{sec:hw}, but with $n_{\mathrm{rounds}}$ rounds.
\item A bare-qubit baseline: a single qubit prepared in $|0\rangle$, delayed by the same compiled duration ($\Delta t$ in backend $\mathrm{dt}$ units) as the Steane circuit, then measured in $Z$. Dynamical decoupling is applied during the delay.
\end{enumerate}
All 8 circuits were submitted as one SamplerV2 batch (job \texttt{d81k4mfoha1c73bkgjlg}, 4096 shots each).

\subsection{Results}

\begin{center}
\begin{tabular}{rrrrr}
\toprule
$n_{\mathrm{rounds}}$ & ISA depth (Steane) & $P(\hat L = 0)$ & $P_{\mathrm{bare}}(Z = 0)$ & $\Delta = \mathrm{logical} - \mathrm{bare}$ \\
\midrule
1 & 229  & 0.8108 & 0.8274 & $-0.017$ \\
2 & 343  & 0.7175 & 0.7637 & $-0.046$ \\
4 & 580  & 0.6492 & 0.6760 & $-0.027$ \\
8 & 1055 & 0.5325 & 0.6235 & $-0.091$ \\
\bottomrule
\end{tabular}
\end{center}

\subsection{Exponential-decay fit and lifetime extraction}

Fitting the model
\[
P(Z = 0; n) - 0.5 \;=\; (P_0 - 0.5)\, e^{-n / \tau_n}
\]
to each curve gives:

\begin{center}
\begin{tabular}{lcc}
\toprule
Curve & $P_0$ & $\tau_n$ (rounds) \\
\midrule
Logical $\hat L = 0$ & $0.906 \pm 0.012$ & $3.57 \pm 0.20$ \\
Bare-qubit Z = 0 & $0.875 \pm 0.008$ & $6.17 \pm 0.34$ \\
\bottomrule
\end{tabular}
\end{center}

The per-round duration (estimated from $(\Delta t_{n=8} - \Delta t_{n=1})/7$) is $\approx 11.80\,\mu\mathrm{s}$. Converting:
\[
\boxed{T_{1,L} = \tau_{n,L} \cdot \Delta t_{\mathrm{round}} = 42.15\,\mu\mathrm{s}, \quad T_{1,\mathrm{phys}} = 72.81\,\mu\mathrm{s}, \quad \frac{T_{1,L}}{T_{1,\mathrm{phys}}} = 0.579.}
\]

The encoded logical qubit retains $58\%$ of the bare-qubit effective lifetime --- a clean engineering number for the cost of $d = 3$ Steane encoding on Heron-r2 at current gate fidelities. For comparison, Heron-r2's published single-qubit $T_1$ is $\approx 280\,\mu\mathrm{s}$; dynamical decoupling on the bare qubit gives $T_2$-eff $\approx 73\,\mu\mathrm{s}$; the encoding then drops this to $T_{1,L} \approx 42\,\mu\mathrm{s}$.

\subsection{Interpretation: we are sub-threshold at $d = 3$}

The logical recovery rate is \emph{below} the bare-qubit baseline at every $n_{\mathrm{rounds}}$ tested, with the gap growing from $-0.017$ at $n_{\mathrm{rounds}} = 1$ to $-0.091$ at $n_{\mathrm{rounds}} = 8$. This means our $d = 3$ Steane code is currently \textbf{below the QEC break-even point on Heron-r2}: a bare qubit with dynamical decoupling preserves $Z$ information better than the encoded logical qubit at matched wall-clock time. The encoding + syndrome-extraction overhead introduces more errors than the $d = 3$ code can correct.

This is the standard expected behaviour for low-distance QEC on near-term superconducting hardware. The published surface-code threshold $p_{\mathrm{th}} \approx 0.7\%$ (per-gate) is achievable only at distance $d \ge 5$ in current architectures; Google Quantum AI's distance-3-vs-5 comparison \cite{willow_precursor} and the Willow distance-7 result \cite{willow} explicitly identify $d = 3$ as a sub-threshold regime even at gate error rates $\sim 0.2$\%. Our \texttt{ibm\_kingston} typical $\mathrm{cz}_{\mathrm{err}} \approx 2 \times 10^{-3}$ is in the same regime.

What our data demonstrates:
\begin{itemize}[leftmargin=*,itemsep=0pt]
\item The Steane code \textbf{runs end-to-end} on Heron-r2 through at least 8 syndrome rounds at ISA depth $\sim 1055$ without circuit-level failures.
\item The logical decay vs rounds is smooth and monotone (logical $Z = 0$ rate drops from 0.81 at $n = 1$ to 0.53 at $n = 8$).
\item The bare-baseline decay is also smooth (0.83 $\to$ 0.62 over the same range), with the bare qubit holding $\approx 9$ percentage points more at $n = 8$.
\item The catcher would correctly flag this as ``smooth'' (no qualitative discontinuity, just monotone decay). A derivative-catcher pass on $P(\hat L)$ vs $n$ identifies $n_{\mathrm{rounds}} \approx 6$ as the steepest-decay region.
\end{itemize}

\textbf{Conclusion: our Steane experiment is a clean baseline below break-even, not a SOTA QEC result.} The natural next step to cross break-even on Heron-r2 is to move to a $d \ge 5$ surface or color code, or to switch to a higher-rate code (Steane, BB) and rely on the lower-overhead/higher-yield trade-off.

\section{Open extensions}\label{sec:extensions}

\begin{enumerate}[leftmargin=*,itemsep=0pt]
\item \textbf{Higher distance}: run distance-5 rotated surface code on a Heron-r2 patch with proper MWPM decoding; this is the smallest distance at which we expect break-even on \texttt{ibm\_kingston} parameters.
\item \textbf{Distance-5 or distance-7 surface code} on a Heron-r2 subgraph using the heavy-hex topology, with proper minimum-weight matching decoding.
\item \textbf{Logical operations}: $X_L$, $Z_L$, $H_L$, and a transversal $\mathrm{CNOT}_L$ between two encoded logical qubits in parallel on the same chip (Heron-r2 has 156 qubits, so 8+ parallel Steane patches fit comfortably with margin).
\item \textbf{Bivariate-bicycle code}: implement a smaller variant of IBM's \texttt{[[72,12,6]]} code on Heron-r2 and compare to Steane.
\end{enumerate}

\section{Reproducibility}

All code, raw counts, and run records are at \href{https://github.com/Zynerji/systrophe}{\texttt{github.com/Zynerji/systrophe}}:

\begin{itemize}[leftmargin=*,itemsep=0pt]
\item Encoder + decoder + syndrome circuit: \texttt{experiments/steane\_logical\_qubit.py}
\item Submission record: \texttt{experiments/results/steane\_kingston\_n4\_submitted.json}
\item Analysis JSON (logical rate, physical rates, sigma): \texttt{experiments/results/steane\_kingston\_n4\_analysis.json}
\item Press-release retraction: \texttt{launch/viral\_press\_release\_156q\_qec.md}
\end{itemize}

\begin{thebibliography}{9}

\bibitem{willow} Google Quantum AI. \emph{Quantum error correction below the surface code threshold}. Nature \textbf{636}, 456 (Dec 2024).
\bibitem{willow_precursor} Google Quantum AI. \emph{Suppressing quantum errors by scaling a surface code logical qubit}. Nature \textbf{614}, 676 (2023).
\bibitem{bb_paper} S.~Bravyi, A.~W.~Cross, J.~M.~Gambetta, D.~Maslov, P.~Rall, T.~J.~Yoder. \emph{High-threshold and low-overhead fault-tolerant quantum memory}. Nature \textbf{627}, 778 (2024).
\bibitem{steane1996} A.~M.~Steane. \emph{Multiple-particle interference and quantum error correction}. Proc.\ R.\ Soc.\ Lond.\ A \textbf{452}, 2551 (1996).
\bibitem{nielsen_chuang} M.~A.~Nielsen, I.~L.~Chuang. \emph{Quantum Computation and Quantum Information}, 10th anniversary ed., Cambridge University Press (2010).

\end{thebibliography}

\end{document}
